\section{Prompt Templates}
\label{appendix:prompts}

For all LLM or MLLM calls, the fallback system prompt (if none specified) is as follows.
\begin{verbatim}
You are a helpful assistant that follows closely the following \
  instructions provided by the user.
\end{verbatim}

In image captioning, we provide captioning instructions and guidelines, context from the source document,
and the image itself to the MLLM.
The image is not part of the textual prompt.
\begin{verbatim}
Based on the surrounding text context and the image content, \
  provide a concise, descriptive caption for this image.

Surrounding context:
{context}

Your task is to produce a caption for the image with context \
  from the original source.
Describe what you see in the image, and give a comprehensive \
  caption that explains the image content as well as the context of the document.
However, do not give information that is too excessive in \
  the caption. 
Remain factual in writing the caption.

The caption should:
1. Describe the visual content clearly
2. Incorporate relevant information from the context
3. Be concise (1-2 sentences)
4. Start with "Image: "
5. NOT express that you based it off the provided context

Caption:
\end{verbatim}

When the knowledge graph option is chosen, we use an LLM to extract entities and relationships from a chunk of text.
The entities and relationships form nodes and edges in the knowledge graph respectively.
We instruct the LLM to extract entities and relationships from the provided chunk, 
and output in the format of a JSON dictionary for easy handling and processing.
\begin{verbatim}
You are an intelligent document processor, and you will assist \
  in storing document information efficiently.
The information below is to be parsed into entities and \
  relationships between them.

# Provided Information
{text}

# Instructions
Read the information and extract entities and relationships \
  between them from the text.

## Entities
For entities you find in the text, provide its **name** and \
  **a short description** of what it is.

## Relationships
For relationships between two entities you find in the text, \
  provide the **source entity**, **target entity**,
the **relationship** between them (relationship type), and \
  their **a short description**.

# Output Format
Give your output in a JSON string similar to the following.
Do **NOT** include any markdown backticks in your output.
Remember that your output will be directly parsed into \
  a Python object.

## Example Output
{{
  "entities": [
    {{"name": "Entity1", "description": "Description of \
                                                entity1"}},
    {{"name": "Entity2", "description": "Description of \
                                                entity2"}}
  ],
  "relationships": [
    {{
      "source": "Entity1", 
      "target": "Entity2", 
      "relationship": "influences", 
      "description": "How entity1 influences entity2"
    }}
  ]
}}

Response:
\end{verbatim}

Communities of nodes in the knowledge graph are detected using Louvain's algorithm.
For each community, we generate a community summary using an LLM.
We include instructions for summarizing and the community's information in the prompt.
\begin{verbatim}
You are a helpful assistant that will assist with \
  data handling tasks.
Previously some documents were processed, and entities \
  and their relationships were extracted to form a \
  knowledge graph.
Some communities were found in the knowledge graph, \
  and your help is needed in organizing these communities.

Summarize the following community of related entities \
  in about {self.config.community_summary_length} words.

Main entities:
{", ".join(f"{node["name"]}" for node in \
                        community_info["nodes"])}

Key relationships:
{"; ".join(f"{relation["source"]} {relation["relationship"]} \
   {relation["target"]}" for relation in \
                          community_info["relationships"])}

Provide a concise summary that captures the main themes 
  and relationships within this community.
\end{verbatim}

We use the Hypothetical Document Embeddings (HyDE) technique for text-to-text retrieval.
The hypothetical document is a zero-shot generated answer to the input query.
While the zero-shot answer may be highly incorrect and hallucinated, 
the generated text contains similar nuances and text patterns similar to retrieval targets that contain evidence to answer the question.
This improves retrieval performance.
We prompt an LLM with short and concise instructions, and the input query.
\begin{verbatim}
Based on the query, generate a detailed hypothetical \
  answer document (approximately 200-300 words).

Query: {query}

Hypothetical Document:
\end{verbatim}

The prompt for final answer generation contains retrieved context and the user's query,
as well as instructions on the style and formation of the answer.
If there are query images, or images are retrieved to be used as context,
the textual prompt to provide such instructions as well.
\begin{verbatim}
You are a helpful assistant that answers questions \
  based on retrieved context in possibly both text and visual form.

{"" if not query_img else f"The user has provided an \
    image in their query, and it is the first \
    {len(query_img)} images among all images \
    presented to you."}
{"" if not context_images else f"The last \
    {len(context_images)} presented to you are \
    relevant visual context to assist you in \
    answering the question."}   

Based on the following context, please answer \
  the question.

Context:
{context}

Question: {query}

Please provide a comprehensive answer.
Always trust the context over your own knowledge.
If you find the context insufficient for answering \
  the question, gently reject answering the question.

You act as a chatbot and are having a conversation \
  with the user.
Provide a natural response as in a conversation; \
  do NOT mention anything about the context in your answer.

Answer:
\end{verbatim}

In testing and benchmarking, we utilize the LLM-as-judge technique \cite{gu2025surveyllmasajudge} to judge final answer correctness.
The prompt to judge generated outputs is
\begin{verbatim}
You are an evaluation judge. Determine if the generated \
  answer is correct given the ground truth answer(s).

Question: {question}

Generated Answer: {generated_ans}

Ground Truth Answer(s): {ground_truths_str}

The generated answer should be considered CORRECT if it:
1. Contains the same factual information as any of the \
    ground truth answers
2. Is semantically equivalent (e.g., "December 1972" \
    matches "14 December 1972 UTC")
3. Minor formatting differences are acceptable \
    (e.g., "James I" vs "James I.")

The generated answer should be considered INCORRECT if:
1. It contains wrong factual information
2. It's too vague or ambiguous
3. It contradicts the ground truth

Respond with ONLY "CORRECT" or "INCORRECT" followed \
  by a brief explanation.

Format:
CORRECT/INCORRECT: <brief explanation>
\end{verbatim}